% !TEX program = xelatex
\documentclass[11pt,a4paper]{article}
\usepackage[a4paper,margin=18mm,headheight=15pt]{geometry}
\usepackage{fontspec}\usepackage{polyglossia}
\setdefaultlanguage{russian}\setotherlanguage{english}
\setmainfont{Arial}\setsansfont{Arial}\setmonofont{Arial}
\usepackage{microtype,enumitem,array,booktabs,tabularx,xcolor,hyperref,fancyhdr,titlesec,tcolorbox,amssymb,lastpage}
\definecolor{MaxPurple}{HTML}{641CF3}\definecolor{MaxBlue}{HTML}{0A84FF}
\definecolor{Ink}{HTML}{16181D}\definecolor{Muted}{HTML}{5E6673}
\definecolor{Panel}{HTML}{F3F6FA}\definecolor{Line}{HTML}{DDE3EC}\definecolor{Alert}{HTML}{D92D20}
\hypersetup{colorlinks=true,linkcolor=MaxPurple,urlcolor=MaxBlue,pdftitle={Спринт 1 - проблема, UX и фундамент},pdfauthor={Команда QR-посещаемости в MAX}}
\pagestyle{fancy}\fancyhf{}
\fancyhead[L]{\small\color{Muted}QR-посещаемость в MAX}
\fancyhead[R]{\small\color{Muted}Спринт 1 \textbullet{} 22-23 сентября}
\fancyfoot[L]{\small\color{Muted}План команды из 4 человек}
\fancyfoot[R]{\small\color{Muted}\thepage\ / \pageref{LastPage}}
\renewcommand{\headrulewidth}{0.4pt}
\titleformat{\section}{\Large\bfseries\color{Ink}}{\colorbox{MaxPurple}{\textcolor{white}{\strut\thesection}}}{0.65em}{}
\titleformat{\subsection}{\large\bfseries\color{Ink}}{\thesubsection}{0.55em}{}
\titlespacing*{\section}{0pt}{1.9ex}{0.9ex}\titlespacing*{\subsection}{0pt}{1.3ex}{0.5ex}
\newlist{tasklist}{itemize}{1}\setlist[tasklist]{label=$\square$,leftmargin=1.6em,itemsep=0.22em,topsep=0.3em}
\setlist[itemize]{leftmargin=1.55em,itemsep=0.2em,topsep=0.3em}
\setlength{\parindent}{0pt}
\newcommand{\code}[1]{\texttt{#1}}
\newcommand{\owner}[1]{\textcolor{MaxPurple}{\textbf{Владелец: #1}}}
\newtcolorbox{goalbox}[1][]{colback=Panel,colframe=MaxPurple,boxrule=0.9pt,arc=2mm,left=3mm,right=3mm,top=2.5mm,bottom=2.5mm,#1}
\newtcolorbox{alertbox}[1][]{colback=white,colframe=Alert,boxrule=0.9pt,arc=2mm,left=3mm,right=3mm,top=2.5mm,bottom=2.5mm,#1}
\begin{document}
\begin{titlepage}\pagecolor{MaxPurple}\color{white}
\vspace*{17mm}{\fontsize{15}{18}\selectfont ПРОЕКТ \textbullet{} QR-ПОСЕЩАЕМОСТЬ В MAX}\par
\vspace{19mm}{\fontsize{36}{42}\selectfont\bfseries Спринт 1\\Проблема, UX\\и технический фундамент}\par
\vspace{9mm}{\Large 22-23 сентября 2026}\par
\vfill
\begin{tcolorbox}[colback=white,colframe=white,arc=3mm,boxrule=0pt]\color{Ink}
\textbf{Цель:} за два дня превратить идею в согласованный MVP и доказать техническую связку MAX Mini App $\rightarrow$ frontend $\rightarrow$ backend $\rightarrow$ PostgreSQL.
\end{tcolorbox}
\vspace{8mm}{\large Команда: Данил, Рома, Никита, Амир}\par
\end{titlepage}\nopagecolor\color{Ink}

\section{Результат спринта}
\begin{goalbox}
К вечеру 23 сентября команда должна иметь подтверждённую проблему, выбранный главный сценарий, карту экранов, низкодетальные макеты, зафиксированную архитектуру и работающий технический скелет. Mini App открывается в MAX, frontend обращается к backend, backend подключён к базе и проверяет подлинность \code{initData}.
\end{goalbox}

\subsection{Границы MVP}
\begin{itemize}
  \item Преподаватель создаёт группу и запускает занятие.
  \item Система показывает динамический QR-код с коротким сроком действия.
  \item Студент открывает ссылку в MAX и отмечает присутствие.
  \item Сервер проверяет пользователя, группу, сессию и токен.
  \item Преподаватель видит список и закрывает сессию.
\end{itemize}
\begin{alertbox}
\textbf{Не входит в этот спринт:} аналитика, интеграция с университетскими системами, роли администратора, сложные отчёты и любые P2-функции. Приоритет - доказать основной путь и снять технические риски MAX.
\end{alertbox}

\section{Распределение ответственности}
\begin{tabularx}{\textwidth}{>{\bfseries}p{25mm}p{34mm}X}
\toprule
Участник & Основная роль & Результат к концу спринта\\
\midrule
Данил & управление, QA, документация & scope, доска, критерии готовности, требования сдачи\\
Рома & product, UX, backend-support & исследование проблемы, As Is / To Be, ценность и метрики\\
Никита & UI/UX, frontend, MAX & карта экранов, макеты, Mini App и клиентская интеграция\\
Амир & backend, security, DevOps & API-скелет, модель данных, Docker, проверка \code{initData}\\
\bottomrule
\end{tabularx}

\section{Задачи Данила}
\owner{Данил}
\begin{tasklist}
  \item Создать доску задач \code{Backlog / Ready / In progress / Review / Done}.
  \item Назначить владельца каждой задачи.
  \item Утвердить рабочий график команды.
  \item Зафиксировать безопасный внутренний дедлайн: 29 сентября, 18:00 МСК.
  \item Определить основной пользовательский сценарий.
  \item Утвердить Must Have / Should Have / Won't Have.
  \item Сформулировать критерии готовности MVP.
  \item Создать Git-репозиторий и правила веток.
  \item Проверить требования к сдаче: код, Docker, README, PDF-презентация, ссылки и тестовые данные.
  \item Получить токен MAX.
  \item Подтвердить окончательный ник бота до создания.
  \item Подготовить шаблон README.
  \item Организовать интервью с преподавателями.
\end{tasklist}

\section{Задачи Ромы}
\owner{Рома - product lead и UX-support}
\begin{tasklist}
  \item Провести или подготовить интервью с преподавателями.
  \item Провести короткие интервью со студентами.
  \item Зафиксировать целевую аудиторию и основной сегмент.
  \item Описать проблему по схеме: пользователь, ситуация, задача, барьер, последствие.
  \item Подготовить As Is.
  \item Подготовить To Be.
  \item Сформулировать ценность для преподавателя.
  \item Сформулировать продуктовую гипотезу.
  \item Выбрать метрики MVP.
  \item Собрать проверяемые источники.
  \item Описать ограничения существующего процесса.
  \item Отделить подтверждённые факты от гипотез команды.
  \item Помочь Никите с пользовательскими путями преподавателя и студента.
\end{tasklist}

\section{Задачи Никиты}
\owner{Никита - UI/UX и frontend lead}
\begin{tasklist}
  \item Создать проект на React, TypeScript и Vite.
  \item Подключить MAX UI.
  \item Подключить MAX Bridge.
  \item Проверить получение \code{initData}.
  \item Проверить получение \code{start\_param}.
  \item Сделать тестовый deep link.
  \item Создать карту экранов преподавателя.
  \item Создать карту экранов студента.
  \item Подготовить low-fidelity макеты.
  \item Сделать экраны входа, выбора группы, QR и результата.
  \item Согласовать состояния экранов и API-контракты с Амиром.
  \item Сделать экран, который открывается в MAX и вызывает backend healthcheck.
\end{tasklist}

\section{Задачи Амира}
\owner{Амир - backend, security и DevOps lead}
\begin{tasklist}
  \item Создать FastAPI-проект.
  \item Подключить PostgreSQL.
  \item Настроить SQLAlchemy и Alembic.
  \item Создать Docker Compose и базовые Dockerfile.
  \item Подготовить \code{.dockerignore} и \code{.env.example} без секретов.
  \item Спроектировать модель данных.
  \item Спроектировать API.
  \item Реализовать healthcheck.
  \item Подготовить серверную проверку MAX \code{initData}.
  \item Реализовать прототип HMAC-SHA256 проверки.
  \item Подготовить начальные миграции.
  \item Создать тестовые данные.
  \item Передать Никите примеры запросов, ответов и кодов ошибок.
\end{tasklist}

\section{Постоянные владельцы направлений}
\begin{tabularx}{\textwidth}{>{\bfseries}p{25mm}X X}\toprule
Участник & Владеет разделами & Поддерживает\\\midrule
Данил & управление, границы MVP, тестирование, документация, презентация, demo, сдача & все направления на приёмке\\
Рома & продуктовое исследование, пилот, масштабирование & MVP, UX, backend, frontend, презентация\\
Никита & UX/UI, MAX integration, frontend & архитектура, backend, документация\\
Амир & архитектура, backend, безопасность, Docker/DevOps & MVP, MAX integration, документация, сдача\\\bottomrule
\end{tabularx}

\section{Ежедневный ритм}
\begin{itemize}
  \item \textbf{10:00 - stand-up, 10 минут:} сделано, план, блокеры.
  \item \textbf{15:00 - быстрая интеграция:} frontend и backend проверяют совместимость.
  \item \textbf{20:00 - вечерняя демонстрация:} показывается только работающий результат.
  \item После демонстрации Данил обновляет приоритеты следующего дня.
\end{itemize}
\begin{alertbox}\textbf{Критический стык: Никита + Амир.} Frontend и backend интегрируются ежедневно; нельзя ждать окончания разработки отдельных веток.\end{alertbox}

\section{Точки интеграции}
\begin{tabularx}{\textwidth}{p{30mm}X>{\raggedright\arraybackslash}p{37mm}}
\toprule
Когда & Совместная проверка & Участники\\
\midrule
22 сен., 15:00 & \code{initData} и \code{start\_param} видны во frontend; согласован auth-контракт & Никита + Амир\\
22 сен., 20:00 & демо Mini App и черновик карты экранов; подтверждена проблема & вся команда\\
23 сен., 15:00 & frontend вызывает backend, backend читает/пишет тестовую БД & Никита + Амир\\
23 сен., 20:00 & sprint review по всем критериям выхода & вся команда\\
\bottomrule
\end{tabularx}

\section{Критерии выхода}
\begin{tasklist}
  \item Подтверждены пользователь, проблема, As Is, To Be и ценность.
  \item Согласованы P0-функции, карты экранов и low-fi макеты.
  \item Зафиксированы стек, модель данных, API-контракты и правила безопасности.
  \item Mini App открывается в мобильном или web-клиенте MAX.
  \item Backend запускается одной Docker-командой, применяет миграции и отвечает на healthcheck.
  \item Frontend успешно вызывает backend.
  \item Сервер проверяет корректный \code{initData} и отклоняет изменённый.
  \item В репозитории нет токенов и секретов.
\end{tasklist}

\section{Рабочий журнал}
\begin{tabularx}{\textwidth}{p{25mm}X X}
\toprule
Дата & Сделано / доказательство & Блокер / решение\\
\midrule
22 сентября & \rule{0pt}{11mm} & \\
23 сентября & \rule{0pt}{11mm} & \\
\bottomrule
\end{tabularx}
\vspace{4mm}
\begin{goalbox}
\textbf{Решение sprint review:} $\square$ принят \quad $\square$ принят с долгом \quad $\square$ не принят.\\[2mm]
\textbf{Незакрытые P0 и владелец:} \rule{105mm}{0.4pt}
\end{goalbox}
\end{document}
